\section{Bối cảnh}

Phần Bối cảnh của nghiên cứu sẽ trình bày bối cảnh về tình trạng thu hồi ấn phẩm khoa học và ý nghĩa của việc theo dõi, phân tích các ấn phẩm bị thu hồi đối với tính toàn vẹn của hoạt động nghiên cứu. Trước hết, phần này khái quát sự phát triển của nghiên cứu trong ngành Khoa học Thông tin và Thư viện và vai trò của các cơ sở dữ liệu thư mục trong việc nhận diện, theo dõi và đánh giá các ấn phẩm bị thu hồi. Tiếp theo, nghiên cứu trình bày các nguyên nhân và đặc điểm của thu hồi ấn phẩm, đồng thời tổng quan các nghiên cứu đã sử dụng phân tích dữ liệu và trắc lượng thư mục để khám phá xu hướng, tác giả, quốc gia, tạp chí, chủ đề và mạng lưới hợp tác liên quan đến các ấn phẩm bị thu hồi. Giai đoạn 2021 - 2026 được lựa chọn nhằm tập trung vào những diễn biến gần đây của hiện tượng thu hồi ấn phẩm, đồng thời bao quát giai đoạn đủ dài để nhận diện xu hướng và biến động theo thời gian; mốc 2021 cần phải chỉ ra có 1 mốc sự kiện gì đặc biệt để bắt đầu theo dõi, trong khi năm 2026 được sử dụng để cập nhật tình hình mới nhất tại thời điểm thu thập dữ liệu. Trên cơ sở đó, phần Background làm rõ khoảng trống nghiên cứu và sự cần thiết của việc kết hợp phân tích dữ liệu với trắc lượng thư mục để cung cấp một bức tranh có hệ thống về các ấn phẩm ngành Khoa học Thông tin và Thư viện bị thu hồi trong giai đoạn nghiên cứu.

Ngoài ra, phần Bối cảnh có thể dẫn dắt đến các câu hỏi nghiên cứu nhằm định hướng cho quá trình phân tích. Các câu hỏi có thể tập trung vào quy mô và xu hướng thu hồi theo thời gian, phân bố theo quốc gia, tổ chức, tạp chí và lĩnh vực/chủ đề nghiên cứu, các nguyên nhân dẫn đến thu hồi, cũng như đặc điểm hợp tác khoa học của các tác giả có ấn phẩm bị thu hồi. Việc xác định các RQ ngay từ phần này giúp liên kết chặt chẽ giữa bối cảnh nghiên cứu, dữ liệu được thu thập, phương pháp trắc lượng thư mục và các kết quả phân tích, đồng thời làm rõ những khía cạnh của hiện tượng retraction mà nghiên cứu hướng tới khám phá. Chẳng hạn, có thể tham khảo ba câu hỏi nghiên cứu sau:\\[3pt]
\highlightlabel{RQ1} Xu hướng số lượng và đặc điểm của các ấn phẩm ngành Khoa học Thông tin và Thư viện bị thu hồi thay đổi như thế nào trong giai đoạn 2021--2026? \\[3pt]
\highlightlabel{RQ2} Những quốc gia, tổ chức, tạp chí và chủ đề nghiên cứu nào có mức độ xuất hiện ấn phẩm bị thu hồi cao nhất? \\[3pt]
\highlightlabel{RQ3} Các nguyên nhân thu hồi và đặc điểm hợp tác khoa học của các ấn phẩm bị thu hồi có những mô hình hoặc xu hướng đáng chú ý nào? \\[3pt]
Các câu hỏi trên có thể tiếp tục được điều chỉnh dựa trên phạm vi dữ liệu, mục tiêu nghiên cứu và các biến số thực tế có trong bộ dữ liệu, qua đó bảo đảm sự nhất quán giữa câu hỏi nghiên cứu, phương pháp phân tích và kết quả nghiên cứu.